\cotefile{85}
\scanpage{30}
conclut le résultat \unsure{analogue} sans restr. sur les
car. résiduelles, car le schéma des \deleted{\unsure{élém.|éléments}} ``réflexions''
dans $\mathrm{GP}(1)_{\mathbb Z}$ est lisse sur $\mathrm{Spec}\,\mathbb Z$. Dans le cas d'
une forme de $\mathrm{Sl}(2)$, OPS que c'est $\deleted{\unsure{\mathrm{Sl}_2|\mathrm{Sl}}}\,(V)$, \deleted{\unsure{donc|i.e.}} \inserted{i.e.}
$u$ est un autom. de $V$ satisfaisant $\deleted{\unsure{u^2=\mathrm{id}}},\ \det u=1$,
\deleted{Mais $u_s\ne\mathrm{id}\quad\forall s\in S$} \unsure{Car|Comme} \deleted{\unsure{donc}} par Hamilton-Cayley
\deleted{\unsure{Tr}}, $u^2-(\mathrm{Tr}\,u)u+1\mathrm{id}=0$, donc $\inserted{u^2=\mathrm{id}\Longleftrightarrow}(\mathrm{Tr}\,u)u=2\mathrm{id}$, \deleted{pour que} \inserted{de cette hyp.}
\deleted{$u^2=\mathrm{id}$, comme $u$ est} il faut conclure ``$\mathrm{Tr}\,u=-2$''. Faisant un calcul infinitésimal,
\deleted{on trouve par erreur que ça coince, d'où \unsure{le}} \deleted{pour en conclure que} \deleted{\unsure{sauf}} aucune solution en
\deleted{nécessité de cette restriction. --}
car. résiduelle 2, OPS qu'on a $\mathrm{Tr}\,u_n=-2$, et \inserted{on va essayer} prouver
$\mathrm{Tr}\,u_{n+1}=\mathrm{Tr}\,u=-2$. \deleted{On aura $\mathrm{Tr}\,u+2=\tau$ avec $\unsure{m\tau}=0$}

% slanted
\margin{\unsure{bien peu|bien au} clair\unsure{…}}

\uline{Proposition} Soit $G$ une forme de $\mathrm{GP}(1)_S$ ou $\mathrm{Sl}(2)_S$, $u$ une
section de $G$, \inserted{$\alpha=A(u)$,} \deleted{telle que $(A(u)+2)(A(u)-2)$ soit inv.,}
$n$ un entier $\geq3$. \deleted{Pour que l'on ait \unsure{Soit $F_n$|Soient $F_n$}} \inserted{On suppose $(\alpha+2)(\alpha-2)$ inv.} Soient

\[
\begin{aligned}
F_n&\in\mathbb Z[U]\\ 
\Psi_n&\in\mathbb Z[U]
\end{aligned}
\]

les deux polynômes \deleted{que} unitaires \deleted{\unsure{sur}} \inserted{à racines deux à deux distinctes} dont les
racines sont 1°) $\zeta+\zeta^{-1}$, où $\zeta$ parcourt les racines
$n$-ièmes \deleted{de} de $1$ \deleted{telles} distinctes de $\pm1$ \enspace 2°) les $\zeta+\zeta^{-1}$,
où $\zeta$ parcourt les racines \uline{primitives} $n$-ièmes de $1$\unsure{.}

\deleted{Pour q}\inserted{Sous ces conditions}

Pour que, $u^n=\mathrm{id}$, il f. et s. que $F_n(\alpha)=1$

Pour que \inserted{on ait} $u^n=\mathrm{id}$ et que $u_s$ soit d'ordre $n$
exactement pour tt $s\in S$, il f. et s. que $\Psi_n(\alpha)=0$.

% slanted
\margin{Sans l'hyp. préalable que $\alpha^2-4$ inv., on peut dire que $[u^n=\mathrm{id}$ et les $u_s$ d'ordre $n$ exactement et $n$ premier aux car. résiduelles$]$ $\Longleftrightarrow$ $[\deleted{\unsure{F}}\Psi_n(\alpha)=0$ \uline{et} $\alpha^2-4$ inv.$]$,}

Principe de démonstration. L'hyp. $\deleted{\alpha^2\ne}\inserted{\alpha\unsure{-|\ne}2}$ inv.
\inserted{i.e. les $u_s$ non unipotents,} implique déjà l'existence d'un unique tore maximal
$T$ qui contient $u$. OPS $T\simeq\mathbb G_m$, donc $u$ donné par une

% archive sheet number: 29
